When vision-language models (VLMs) are fine-tuned for robotic action prediction, which representational changes are associated with training and which are inherited from the pretrained backbone?
We compare trained and reference checkpoints of three Vision-Language-Action (VLA) models across four LIBERO benchmark task suites via layer-wise probes, attention profiling, and counterfactual perturbations.
We find that VLM-to-VLA fine-tuning is associated with a consistent exit-layer collapse in action-probe decodability ($\Delta R^2 = 0.16$--$0.46$) that is absent in VLA-to-VLA transitions ($\Delta R^2 < 0.10$) and in three natural language processing (NLP) instruction-tuning controls ($\Delta R^2 = 0.000$), suggesting this reorganization is specific to action-oriented training rather than a generic fine-tuning artifact.
Complementary analyses reveal that trained VLAs develop substantial sensitivity to instruction content while the dominant attention patterns at early layers largely pre-exist in pretrained backbones.
These findings provide a cross-architecture decomposition of instruction routing in fine-tuned VLAs, disentangling training-associated representational changes from inherited structure.
